% Pursuing Stacks (cotes 134-2 to 134-8): archive cote 134-5, site page 62 (PDF page 63), typescript page 249.
% typescript. Draft compared: ps4.tex:3125-3188.
% AI-transcribed by Claude Opus 5.5 (claude-opus-5-5, adaptive thinking, effort high), Message Batches API, batch ps2 (msgbatch_01NVwr1Yuo8qJf5rykeGBXPK), 2026-11-12; not yet checked by eye.
\scanpage{134-5/62}
\mspage{249}

\xout{Condition} The first statement a) can be rephrased,  by saying that
\xout{Mx} the category $M_c$ of contractible objects of $M$ is an \uline{aspherator} .  To
prove this,  we use the fact that $M_c$ is stable under finite products,
which implies that we can find a small full subcategory \xout{of} $C$ of $M_c$ ,
stable under such products,  and which generates the contractibility
structure.  The stability condition upon $C^{\unsure{*}}$ implies that $C^{\wedge}$\hand{/}is totally
aspheric.  Therefore a) and b) will be proved,  if we prove that any
small full subcategory $C$ of $M_c$ ,  such that $C^{\wedge}$ is totally aspheric,
satisfies the conditions of prop.1 (with $N=M_c$) ,  namely that for any
$x$ in $M_c$ ,  $i^{*}(x)$ is an aspheric object of $C^{\wedge}$,  where $i\colon C\to M$  (or $C\to M_c$,
equivalently) is the inclusion functor.  But from the fact that $i^{*}$ commu-%
tes with finite products,  it follows that $i^{*}(x)$  is contractible in $C^{\wedge}$,
for the \xout{homotopy} homotopy interval structure admitting as a generating
family of homotopy intervals the intervals in $C^{\wedge}$ made up with objects of
$C$ .  The assumption of total asphericity upon $C^{\wedge}$ implies that these
homotopy intervals are aspheric over $e_{C^{\wedge}}$ ,  and from this follows \xout{that}
(as already used a number of times earlier)  that $i^{*}(x)$ \inserted{too} is aspheric
over $e_{C^{\wedge}}$ ,  and hence aspheric as $e_{C^{\wedge}}$ is aspheric (because of the assump-%
tion of total asphericity).  This proves a) and b),  and c) follows via
prop. 5 (p. 239).

We'll call of course the asphericity structure described in prop.2
the \uline{asphericity} \uline{structure} \uline{generated} \uline{by the given contractibility struc\xout{tur}}
\uline{ture}.
\uline{Remarks}\inserted{\unsure{1}} Clearly,  not every asphericity structure can be generated by a
contractibility structure,  as a necessary condition (presumably not a
sufficient one) is total asphericity.  We don't expect either,  in case
\margin{X} it can,  that the \inserted{generating} contractibility structure is uniquely defined;  howe-%
ver,  we do expect in this case that there should exist a canonical
(largest) choice - we\over{ll}{\unsure{'ll}} have to come back upon this in due course. For
the time being,  let's only remark that all modelizing asphericity struc-%
tures met \inserted{with} so far,  it seems,  do come from contractibility structures.

It occurs to me that \xout{in} the last statement  was a little hasty  -
after all we have met with test categories which are not strict ones,
hence ($A^{\wedge}$ , $A^{\wedge}_{as}$ )  is a modelizing asphericity structure  (it would be
enough even that $A$ be a pseudo-test category),  which is'nt totally asphe
ric,  and a fortiori does not come from a contractibility structure. \hand{Thus}

\margin{62}
